\section{Introduction}

\subsection{Background \& Motivation}

Over the past century, computing capabilities have grown at an unprecedented rate. The
\textbf{Apollo Guidance Computer} (AGC), a pioneering system developed for the Apollo
program, operated with just a \textit{2 MHz CPU}, \textit{4 KB of RAM}, and \textit{72
KB of ROM}. Despite its severe hardware constraints, the AGC’s operating system
introduced a series of innovations that remain remarkably relevant: a hybrid task
scheduler combining cooperative and preemptive mechanisms, robust error detection with
system-wide recovery, and efficient use of limited memory through careful task
prioritization and synchronization.

One of the AGC's most groundbreaking design philosophies was its embrace of
\textit{restart-driven fault tolerance}. Rather than attempt complex error handling
routines in hardware-constrained conditions, the AGC favored quick, controlled restarts
to recover from faults, enabling resilience even under critical mission conditions—as
dramatically demonstrated during the Apollo 11 lunar descent when 1201 and 1202 alarms
threatened to abort the landing.

Today, a device like the \texttt{micro:bit}—costing less than \pounds20—offers
significantly greater raw computing power than the AGC. Yet in embedded systems, the
fundamental challenges remain: \textit{real-time responsiveness}, \textit{reliability},
\textit{fault tolerance}, and \textit{resource efficiency} are still paramount,
especially for mission- or safety-critical applications.

This project draws inspiration not only from the AGC's design but also from the modular
robustness of modern Unix systems like Linux, and from the elegant abstraction
facilities provided by the C++ Standard Library. The goal is to reimagine the AGC’s
philosophies for a modern embedded context: combining historical pragmatism with
contemporary programming paradigms.

\subsection{Project Summary}

This work presents \textbf{Apollo-RTOS}: a fully-featured, real-time operating system
for the \texttt{micro:bit} v1, deeply inspired by the AGC and informed by the design
philosophies of Linux and C++20. Apollo-RTOS adapts the AGC’s core concepts—most
notably its hybrid cooperative-preemptive scheduling model and restart-driven recovery
mechanism—while integrating modern operating system abstractions essential for
contemporary embedded platforms.

Key architectural features include: 
\begin{itemize}[left=10pt] 
    \item A lightweight, deterministic \textbf{hybrid scheduler} capable of supporting
        up to 16 concurrent processes, with synchronization primitives ensuring
        race-free concurrent execution. [\autoref{sec:sched}] 
    \item A \textbf{fault-tolerant recovery system} based on restart groups, enabling
        processes to persist across resets or power failures. [\autoref{sec:rec}] 
    \item A modular \textbf{file system} that unifies flash and FRAM storage, providing
        durable, high-endurance storage critical for recovery and data integrity.
        [\autoref{sec:fs}] 
    \item A \textbf{Unix-like shell interface} for process management, file
        manipulation, and system interaction, offering a familiar development
        environment. [\autoref{sec:shell}] 
    \item An \textbf{automated testing framework} and build infrastructure to validate
        components, accelerate development, and enable systematic debugging.
        [\autoref{sec:testing}] 
    \item A \textbf{C++20-based implementation} leveraging RAII, templates, constexpr
        programming, and static typing, inspired by the robustness and clarity of the
        C++ Standard Library. 
\end{itemize}

Apollo-RTOS is built around a \textit{monolithic kernel architecture}, balancing
simplicity and performance. Critical system services—including drivers for UART, I2C,
and display peripherals—are designed as lightweight, reusable modules. File system
operations are abstracted over a block device layer, supporting both high-endurance
FRAM and flash memory with transparent access semantics.

Beyond reinterpreting the AGC’s innovations, Apollo-RTOS demonstrates how combining
minimalist principles with modern software engineering can produce resilient,
extensible, and maintainable real-time operating systems. It offers a practical
blueprint for building robust, modular systems under tight resource constraints—an
increasingly valuable skill in the landscape of embedded and edge computing.

The complete source code for Apollo-RTOS is available at:
\[
   \impeq{\text{\url{https://github.com/AnglyPascal/apollo-rtos}}}
\]

\lstinputlisting[style=mytty, caption={Landing screen of Apollo-RTOS on
\texttt{minicom}}]{code/ex1.sh}

% \begin{minipage}{.5\linewidth}
%     \includegraphics[width=.9\linewidth]{ex1}
% \end{minipage}\hfill%
% \begin{minipage}{.5\linewidth}
%     \includegraphics[width=.9\linewidth]{ex2}
% \end{minipage}
